\chapter{Perfil Profissional do Egresso}

A formação do Engenheiro de Computação envolve uma base sólida em matemática, ciência da computação e engenharia, com ênfase em áreas como eletrônica, sistemas embarcados, controle e automação, telecomunicações e redes de computadores. Além da competência técnica, o curso promove o desenvolvimento de habilidades interpessoais, de liderança, comunicação, trabalho em equipe e pensamento crítico, associadas a uma consciência ética, social e ambiental.

Em consonância com o CC2020, espera-se que o egresso deste curso:

\begin{itemize}
    \item Atue com autonomia intelectual, curiosidade investigativa e compromisso com a aprendizagem ao longo da vida;
    \item Reconheça e se comprometa com a diversidade, a inclusão, a equidade e a sustentabilidade em suas práticas profissionais;
    \item Compreenda os impactos sociais, éticos e ambientais dos sistemas computacionais na vida das pessoas e na sociedade;
    \item Domine os fundamentos técnicos para analisar, projetar, implementar e manter soluções computacionais e tecnológicas robustas, seguras, eficientes e escaláveis;
    \item Comunique-se com clareza e eficiência com públicos técnicos e não técnicos, em português e em inglês;
    \item Seja capaz de atuar de forma responsável e proativa em contextos colaborativos, interdisciplinares, multiculturais e globais;
    \item Exerça liderança técnica e empreendedora na condução de projetos e iniciativas inovadoras.
\end{itemize}

Em acordo com a Resolução CNE/CES nº 5/2016, são competências esperadas do egresso de \nomecursonovo:

\begin{itemize}
    \item Planejar, projetar, implementar, testar, verificar e validar sistemas computacionais digitais, incluindo computadores, sistemas embarcados, redes, plataformas distribuídas e sistemas de automação;
    \item Compreender e aplicar normas e boas práticas de segurança em sistemas computacionais;
    \item Desenvolver processadores específicos, arquiteturas otimizadas e software para sistemas embarcados e de tempo real;
    \item Projetar e administrar redes de computadores e infraestruturas de comunicação;
    \item Realizar estudos de viabilidade técnico-econômica e avaliar impactos tecnológicos, financeiros, ambientais e sociais de soluções computacionais;
    \item Conhecer e respeitar os direitos de propriedade intelectual e os marcos legais e regulatórios aplicáveis à computação e à engenharia.
\end{itemize}

Conforme o CONFEA~\cite{CONFEA:473:2002}, o título profissional conferido é o de \textbf{Engenheiro de Computação (Eng. Comp.)}, com atribuições regulamentadas pelo Sistema CONFEA/CREA. O egresso deste curso estará habilitado para atuar nas mais diversas áreas da engenharia computacional, como:

\begin{itemize}
    \item Sistemas embarcados e automação industrial;
    \item Arquiteturas de computadores e sistemas digitais;
    \item Redes de computadores e telecomunicações;
    \item Cibersegurança e sistemas críticos;
    \item Internet das Coisas (IoT) e sistemas inteligentes;
    \item Inteligência artificial e aprendizado de máquina;
    \item Engenharia de software e integração hardware-software;
    \item Computação de alto desempenho e em nuvem.
\end{itemize}

O perfil do egresso reflete o compromisso do curso com uma formação orientada por competências, fundamentada em valores de excelência, responsabilidade ética, cidadania e transformação social.

\section{Egresso do Curso de \nomecursonovo}

O egresso do curso de \nomecursonovo da Universidade Federal de Uberlândia é um(a) profissional com sólida formação científica, tecnológica e humanística, preparado(a) para atuar no projeto, desenvolvimento, integração, gerenciamento e inovação de sistemas computacionais e eletrônicos. Sua formação o(a) capacita a unir hardware e software de forma crítica, criativa, ética e responsável, com alta capacidade de adaptação a cenários multidisciplinares e em constante transformação.

A formação ofertada permite ao egresso desenvolver competências integradas nos seguintes eixos:

\begin{enumerate}
\item Fundamentos Técnico-Científicos

\begin{itemize}
    \item Dominar os conhecimentos fundamentais de matemática, ciência da computação, engenharia elétrica e eletrônica, integrando teoria e prática para o projeto e análise de sistemas computacionais;
    \item Planejar, especificar, implementar, testar, validar e manter sistemas digitais, embarcados, de automação, redes e dispositivos inteligentes, conforme princípios da engenharia;
    \item Aplicar o pensamento computacional e a abstração algorítmica para resolver problemas complexos com eficiência, qualidade e segurança;
    \item Analisar viabilidade técnica e econômica de soluções computacionais, considerando requisitos de desempenho, custo, escalabilidade e sustentabilidade.
\end{itemize}
    
\item Inovação, Criatividade e Empreendedorismo

\begin{itemize}
    \item Desenvolver produtos, serviços e soluções tecnológicas inovadoras, integrando criatividade com conhecimento técnico;
    \item Compreender e aplicar direitos de propriedade intelectual, marcos legais e regulatórios pertinentes à engenharia de computação;
    \item Atuar de forma proativa em ambientes empreendedores, reconhecendo oportunidades de negócio e liderando projetos de base tecnológica.
\end{itemize}
    
\item Ética, Sustentabilidade e Impacto Social

\begin{itemize}
    \item Refletir criticamente sobre o impacto social, ambiental, ético e cultural das tecnologias da computação na vida cotidiana e na sociedade;
    \item Projetar soluções que respeitem direitos humanos, diversidade, acessibilidade e inclusão;
    \item Agir com responsabilidade, autonomia e consciência social, contribuindo para o desenvolvimento sustentável e equitativo.
\end{itemize}
    
\item Comunicação, Colaboração e Liderança

\begin{itemize}
    \item Comunicar-se com clareza, precisão e adequação com diferentes públicos, de forma oral e escrita, em português e inglês técnico;
    \item Trabalhar eficazmente em equipes interdisciplinares, multiculturais e distribuídas, assumindo papéis de liderança técnica e facilitando a cooperação;
    \item Demonstrar empatia, responsabilidade compartilhada e respeito à diversidade no trabalho em grupo e na mediação de conflitos .
\end{itemize}
    
\item Aprendizagem Contínua e Adaptabilidade
    
\begin{itemize}
    \item Demonstrar autonomia na gestão do próprio aprendizado, reconhecendo a necessidade de atualização constante diante da evolução científica e tecnológica;
    \item Adaptar-se a novos contextos, ferramentas, metodologias e ambientes de trabalho, mantendo curiosidade e resiliência frente à mudança;
    \item Avaliar criticamente tecnologias emergentes e aplicar princípios de inovação de forma ética e eficaz.
\end{itemize}
    
\item Atuação Profissional e Global
    
\begin{itemize}
    \item Atuar de forma competente, ética e inovadora nos seguintes campos:
    
    \begin{itemize}
        \item Sistemas embarcados e de tempo real
        \item Engenharia de software e arquiteturas computacionais
        \item Automação industrial e controle
        \item Redes de computadores e telecomunicações
        \item Sistemas distribuídos e computação em nuvem
        \item Cibersegurança e confiabilidade
        \item Inteligência artificial e ciência de dados
    \end{itemize}
    
    \item Intervir com competência técnica em ambientes locais, nacionais e internacionais, respeitando legislações, culturas e expectativas sociais.
\end{itemize}
\end{enumerate}

\section{Habilidades e Competências}

O profissional egresso do curso de \nomecursonovo da UFU deve desenvolver um conjunto articulado de conhecimentos técnicos, habilidades práticas e disposições profissionais, de modo a atuar com excelência, responsabilidade social e capacidade de inovação, conforme estabelecem as Diretrizes Curriculares Nacionais da área de Computação~\cite{MEC:CNE:CES:5:2016} e do \textit{Computing Curricula 2020}~\cite{ACM:CC2020}.

\begin{enumerate}[label=\Roman*.]
\item Competências Técnicas~\cite{MEC:CNE:CES:5:2016}

O curso visa assegurar ao egresso as seguintes competências técnicas:

\begin{itemize}
    \item Planejar, projetar, implementar, testar, verificar e validar sistemas computacionais digitais, incluindo computadores, redes, sistemas embarcados e de automação;
    \item Desenvolver software e hardware para sistemas de tempo real e dispositivos específicos;
    \item Projetar e administrar redes de computadores e plataformas de comunicação;
    \item Compreender e aplicar normas de segurança, boas práticas de desenvolvimento e marcos legais da área;
    \item Gerenciar projetos e manter sistemas computacionais complexos;
    \item Realizar análises de viabilidade técnico-econômica e avaliar impactos tecnológicos, sociais e ambientais;
    \item Conhecer e aplicar direitos de propriedade intelectual e princípios éticos no exercício profissional;
    \item Analisar, avaliar e selecionar plataformas de hardware e software adequadas a aplicações embarcadas, distribuídas ou de missão crítica.
\end{itemize}

\item Habilidades Cognitivas e Práticas~\cite{ACM:CC2020}

Complementando as competências normativas, espera-se que o egresso demonstre:

\begin{itemize}
    \item Raciocínio lógico, computacional e analítico para resolver problemas complexos;
    \item Capacidade de abstração, modelagem e decomposição de sistemas;
    \item Comunicação eficaz com públicos técnicos e não técnicos, em diferentes meios e idiomas;
    \item Capacidade de colaboração em equipes interdisciplinares e multiculturais;
    \item Adaptabilidade a novos contextos tecnológicos e metodológicos;
    \item Organização e gestão do tempo e do próprio aprendizado;
    \item Liderança técnica em projetos, com tomada de decisão embasada e eficaz;
    \item Capacidade de avaliação crítica de sistemas computacionais e inovação tecnológica.
\end{itemize}

\item Disposições e Atitudes Profissionais~\cite{ACM:CC2020}

Além dos conhecimentos e habilidades, busca-se que o egresso adote atitudes compatíveis com a atuação ética e transformadora, tais como:

\begin{itemize}
    \item Comprometimento com a aprendizagem contínua e o aprimoramento pessoal e profissional;
    \item Responsabilidade social, ambiental e ética no uso e desenvolvimento de tecnologias;
    \item Respeito à diversidade, à equidade e à inclusão nos ambientes de trabalho e pesquisa;
    \item Curiosidade investigativa, autonomia intelectual e pensamento crítico;
    \item Empatia, escuta ativa e comportamento colaborativo;
    \item Engajamento com o impacto positivo da computação na vida das pessoas e na sociedade.
\end{itemize}
\end{enumerate}

Essa integração entre os referenciais nacionais e internacionais garante que o egresso do curso de \nomecursonovo da UFU esteja preparado não apenas para os desafios técnicos da profissão, mas também para liderar transformações sociais e tecnológicas em nível local e global.

\section{Exercício Profissional no Brasil}

O exercício no Brasil da profissão de engenheiro é regulamentado desde 1966~\cite{Lei:5194:1966}. A verificação e a fiscalização do exercício e atividades da profissão são exercidas por um Conselho Federal de Engenharia, Arquitetura e Agronomia (CONFEA), e por um Conselho Regional de Engenharia e Agronomia (CREA). Os profissionais habilitados na forma estabelecida na referida lei só poderão exercer a profissão após o registro no Conselho Regional sob cuja jurisdição se achar o local de sua atividade.

Após a publicação das Diretrizes Curriculares Nacionais dos cursos da área de Computação, o CREA passou a registrar sem restrições os profissionais de Engenharia de Computação egressos dos cursos em conformidade com a nova legislação.

A Resolução CONFEA nº 473/2002, atualizada em 31 de março de 2017~\cite{CONFEA:473:2002}, determina que, junto ao sistema CONFEA/CREA, o registro de um profissional de Engenharia de Computação seja identificado como:

\begin{itemize}
    \item Grupo: 1 -- Engenharia
    \item Modalidade: 2 -- Eletricista
    \item Nível: 1 -- Graduação
    \item Código: 121-01-00 -- Engenheiro(a) de Computação (Eng. Comp.)
\end{itemize}

Ainda, de acordo com a Resolução CONFEA nº 1.073/2016~\cite{CONFEA:1073:2016}, são designadas as seguintes atividades para o Engenheiro(a) de Computação profissional:

\begin{enumerate}[label=\Roman*.]
    \item Gestão, supervisão, coordenação, orientação técnica.
    \item Coleta de dados, estudo, planejamento, anteprojeto, projeto, detalhamento, dimensionamento e especificação.
    \item Estudo de viabilidade técnico-econômica e ambiental.
    \item Assistência, assessoria, consultoria.
    \item Direção de obra ou serviço técnico.
    \item Vistoria, perícia, inspeção, avaliação, monitoramento, laudo, parecer técnico, auditoria, arbitragem.
    \item Desempenho de cargo ou função técnica.
    \item Treinamento, ensino, pesquisa, desenvolvimento, análise, experimentação, ensaio, divulgação técnica, extensão.
    \item Elaboração de orçamento.
    \item Padronização, mensuração, controle de qualidade.
    \item Execução de obra ou serviço técnico.
    \item Fiscalização de obra ou serviço técnico.
    \item Produção técnica e especializada.
    \item Condução de serviço técnico.
    \item Condução de equipe de produção, fabricação, instalação, montagem, operação, reforma, restauração, reparo ou manutenção.
    \item Execução de produção, fabricação, instalação, montagem, operação, reforma, restauração, reparo ou manutenção.
    \item Operação, manutenção de equipamento ou instalação.
    \item Execução de desenho técnico.
\end{enumerate}

De acordo com a Resolução CONFEA nº 380/1993~\cite{CONFEA:380:1993} e nº 218/1973~\cite{CONFEA:218:1973}, essas atividades, no âmbito da Engenharia de Computação, são referentes exclusivamente a:

\begin{enumerate}[label=\Roman*.]
    \item Materiais elétricos e eletrônicos;
    \item Equipamentos eletrônicos em geral;
    \item Sistemas de comunicação e telecomunicações;
    \item Sistemas de medição e controle elétrico e eletrônico; seus serviços afins e correlatos;
    \item Análise de sistemas computacionais, seus serviços afins e correlatos.
\end{enumerate}

Destaca-se que o curso de \nomecursonovo da FEELT/UFU, em sua versão atual, não se enquadra no requerido pela Resolução CONFEA nº 380/1993~\cite{CONFEA:380:1993} em seu art. 1º, § 2º. Caso seja de interesse do(a) discente incluir as referências do art. 8º da Resolução CONFEA nº 218/1973 (geração, transmissão, distribuição e utilização da energia elétrica; equipamentos, materiais e máquinas elétricas; sistemas de medição e controle elétricos; seus serviços afins e correlatos)~\cite{CONFEA:218:1973}, o mesmo precisará cursar disciplinas extracurriculares e requisitar especificamente para o CREA a extensão de referências em suas atividades.
